\scanpage{329}
\mspage{2}

classe d'homotopie \unsure{du chemin|des chemins} $c$ (i.e. son image dans $\mathrm{Fl}(\Pi(X))$ est égale à celle de $\unsure{u_0|\alpha_0}$. Pour le voir, choisissons pour \unsure{tt|tout} $i\in[0,\unsure{2N|N}]$ une flèche $\alpha_i:x_i\longrightarrow a$ \unsure{.} (ça existe grâce à l'hypothèse b)).

Je dis que pour chaque $i\in[0,2N]$, le triangle

\diagram{\begin{tikzcd}[ampersand replacement=\&]
 x_0\arrow[rr,"c_i"]\arrow[dr,dash,"\alpha_0"'] \& \& x_i\arrow[dl,dash,"\alpha_i"]\\
 \& a \&
\end{tikzcd}}

dans $\Pi(X)$ est commutatif, où $c_i$ désigne la classe d'homotopie du chemin de $x_0$ à $\unsure{x_i|x}$ défini par les $u_i$ ($i\in[1,i]$). Nous le faisons par récurrence sur $i$, \deleted{de sorte qu'on est} \inserted{l'assertion est} triviale si $i=0$ ($c_0$ représente l'application identique.)

\unsure{Le|On} \deleted{\unsure{est|a|un}} diagramme

\diagram{\begin{tikzcd}[ampersand replacement=\&,column sep=small]
 x_0\arrow[r,dash,"c_i"]\arrow[rr,dash,bend left=18,"c_{i+1}"]\arrow[dr,"\alpha_0"'] \& x_i\arrow[r,"u_{i+1}"]\arrow[d,"\alpha_i"] \& x_{i+1}\arrow[dl,"\alpha_{i+1}"]\\
 \& a \&
\end{tikzcd}}

montre qu'il suffit de prouver que pour $i\in[0,2N-1]$ le triangle

\diagram{\begin{tikzcd}[ampersand replacement=\&]
 x_i\arrow[rr,dash,"u_{i+1}"]\arrow[dr,"\alpha_i"'] \& \& x_{i+1}\arrow[dl,"\alpha_{i+1}"]\\
 \& a \&
\end{tikzcd}}

est commutatif dans $\Pi(X)$. (\uline{NB} $u_{i+1}$ est en sens direct ou rétrograde, suivant que $i$ est pair ou impair.) Mais cela résulte du fait que \deleted{\unsure{Hom}} l'image de $\operatorname{Hom}_X(x_{\unsure{i|1}},a)$ (resp $\operatorname{Hom}_{\deleted{\unsure{X}}X}(\unsure{x_{i+1}},\unsure{?|a})$

% archive sheet number: 329
